\usepackage{fontspec}
\setmainfont[BoldFont=texgyretermes-bold.otf,ItalicFont=texgyretermes-italic.otf,BoldItalicFont=texgyretermes-bolditalic.otf]{texgyretermes-regular.otf}
\setsansfont[BoldFont=texgyreheros-bold.otf]{texgyreheros-regular.otf}
\setmonofont[Scale=MatchLowercase]{lmmono10-regular.otf}
\setjamainfont[BoldFont=HaranoAjiMincho-Bold.otf]{HaranoAjiMincho-Regular.otf}
\setjasansfont{HaranoAjiGothic-Medium.otf}
\usepackage{amsmath,amssymb,bm}
\usepackage{graphicx,booktabs,longtable,array,tabularx}
\usepackage{xcolor}
\usepackage{caption}
\usepackage{enumitem}
\usepackage{fancyhdr}
\usepackage{url}
\usepackage[unicode,bookmarksnumbered,colorlinks=true,linkcolor=blue!45!black,citecolor=blue!45!black,urlcolor=blue!45!black]{hyperref}
\hypersetup{pdftitle={音響によるプール沸騰の熱流束回帰と沸騰開始判定におけるモデル統合の効果と限界（修士論文草稿）},pdfauthor={柴崎 陸},pdfsubject={2026年10月19日時点の保存結果に基づく執筆用草稿}}
\graphicspath{{figures/}}
\setlength{\emergencystretch}{3em}
\setlength{\parindent}{1em}
\setlength{\parskip}{0.2em}
\renewcommand{\arraystretch}{1.16}
\captionsetup{font=small,labelfont=bf}
\setlist{itemsep=0.25em,topsep=0.4em}
\pagestyle{fancy}
\fancyhf{}
\fancyhead[L]{\small 修士論文草稿：2026年10月19日}
\fancyhead[R]{\small\nouppercase{\leftmark}}
\fancyfoot[C]{\thepage}
\setlength{\headheight}{17pt}
\renewcommand{\headrulewidth}{0.3pt}
\newcommand{\kw}{\ensuremath{\mathrm{kW/m^2}}}
\newcommand{\wm}{\ensuremath{\mathrm{W/m^2}}}
\newcommand{\mseunit}{\ensuremath{(\mathrm{kW/m^2})^2}}
\newcommand{\qonb}{q_{\mathrm{ONB},d}}
\newcommand{\qhat}{\hat q}
\newcommand{\need}[2]{\begin{quote}\small\color{red!65!black}\textbf{【要確認：#1】} #2\end{quote}}
\newcommand{\finding}[1]{\par\noindent\textbf{#1}\par}
\newcommand{\source}[1]{\par\smallskip{\footnotesize 出典：#1}\par}
\newcolumntype{L}[1]{>{\raggedright\arraybackslash}p{#1}}
\newcolumntype{Y}{>{\raggedright\arraybackslash}X}
\setcounter{tocdepth}{2}
\setcounter{secnumdepth}{2}
